\chapter{ВСТУП}

\section{Контекст та актуальність}

Промислові підприємства стикаються з несподіваними поломками обладнання. Це коштує дорого і забирає час. Інтернет речей в комбінації з машинним навчанням можуть допомогти, адже вони дозволяють відстежувати обладнання в реальному часі та передбачати поломки. Традиційні системи SCADA та корпоративні платформи IoT теж виконують ці функції, але мають недоліки:

\begin{itemize}
    \item висока вартість --- рішення від SAP, Siemens, GE коштують від сотень тисяч до мільйонів доларів, малі та середні підприємства не можуть собі цього дозволити;
    \item складна інтеграція --- впровадження SCADA займає від 6 місяців до 2 років, потрібні спеціалісти для налаштування, відсутність стандартів ускладнює підключення до існуючої інфраструктури;
    \item погана масштабованість --- монолітна архітектура створює вузькі місця при збільшенні кількості сенсорів, горизонтальне масштабування потребує серйозних змін або повної заміни системи;
    \item прості сповіщення --- традиційні системи використовують фіксовані пороги для окремих параметрів, вони не бачать складні аномалії, коли проблема виникає від комбінації параметрів;
    \item ізоляція даних --- коли кілька компаній використовують спільну інфраструктуру, потрібна надійна ізоляція промислових даних, більшість традиційних систем не розраховані на багатокористувацький режим.
\end{itemize}

Ці обмеження створюють потребу в новій платформі, яка має поєднувати доступність технологій з відкритим кодом, корпоративну надійність та машинне навчання для виявлення аномалій.


\section{Проблема}

Основні проблеми промислового Інтернету речей:

\begin{itemize}
    \item високий поріг входу --- підприємствам потрібне надійне відстеження обладнання, але корпоративні рішення закриті та прив'язують до одного постачальника;
    \item технічна складність --- створити власну систему важко, потрібна експертиза в розподілених системах, обробці даних в реальному часі, машинному навчанні та безпеці, це вимагає команду спеціалістів та інвестиції в розробку;
    \item незручні інтерфейси --- інтерфейси SCADA створені інженерами для інженерів, операційним менеджерам та керівникам потрібно більше часу для освоєння цих інструментів.
\end{itemize}


\section{Мета та завдання}

Мета роботи --- створити багатокористувацьку платформу для промислового моніторингу, яка масштабується, працює в реальному часі та виявляє аномалії через машинне навчання на технологіях з відкритим кодом.

Завдання дослідження:

\begin{enumerate}
    \item спроектувати мультитенантну архітектуру з окремою схемою бази даних для кожного користувача та горизонтальним масштабуванням через партиції Kafka та воркери Spark;
    \item реалізувати обробку даних в реальному часі через Apache Kafka та Apache Spark Structured Streaming з агрегаціями за періоди (1хв, 5хв, 1год);
    \item інтегрувати виявлення аномалій через три алгоритми (Isolation Forest, Random Forest, XGBoost) з MLflow для відстеження експериментів та версій моделей;
    \item створити систему сповіщень з трьома режимами детекції (порогові, ML, комбіновані);
    \item розробити REST API на FastAPI з JWT автентифікацією та документацією OpenAPI;
    \item налаштувати моніторинг через Prometheus та Grafana;
    \item протестувати систему на продуктивність, масштабованість та якість ML моделей.
\end{enumerate}


\section{Вимоги до системи}

\subsection{Функціональні вимоги}

\begin{itemize}
    \item Ф-001 Багатокористувацька архітектура (критичний) --- кожна компанія має власну схему в PostgreSQL/TimescaleDB; дані повністю ізольовані; економія 90-95\% місця через стиснення.
    \item Ф-002 Обробка в реальному часі (критичний) --- Apache Kafka та Apache Spark обробляють дані в реальному часі; автоматичні агрегації за 1хв, 5хв, 1год; затримка не більше 5 секунд.
    \item Ф-003 Виявлення аномалій (високий) --- три алгоритми (Isolation Forest, Random Forest, XGBoost); MLOps інфраструктура (MLflow, реєстр моделей, MinIO).
    \item Ф-004 Сповіщення (високий) --- CRUD операції для правил; три режими детекції (порогові, ML, комбіновані); публікація в Kafka.
    \item Ф-005 API (критичний) --- FastAPI шлюз; JWT автентифікація (7 днів); автоматичний search\_path для ізоляції; пул з'єднань; OpenAPI документація.
    \item Ф-006 Ієрархія обладнання (високий) --- гібридна архітектура з нормалізованими таблицями (Компанія → Об'єкт → Обладнання → Сенсор) та денормалізованими вимірюваннями.
    \item Ф-007 Масштабування (середній) --- горизонтальне масштабування через партиції Kafka та воркери Spark без змін коду.
\end{itemize}

\subsection{Нефункціональні вимоги}

\begin{itemize}
    \item НФ-001 Продуктивність (високий) --- асинхронне програмування FastAPI; пул з'єднань; кеш Redis (60с); мікропакетна обробка Spark; неперервні агрегати TimescaleDB.
    \item НФ-002 Безпека (критичний) --- JWT (HS256, 7 днів); Bcrypt (складність 12); ізоляція search\_path; параметризовані SQL; CORS; аудит.
    \item НФ-003 Надійність (критичний) --- доступність 99.99\%; реплікація Kafka $\geq$ 2; контрольні точки Spark; резервні копії; автовідновлення; політики перезапуску Docker.
    \item НФ-004 Зберігання (високий) --- гіпертаблиці з щоденними фрагментами; стиснення після 7 днів (90-95\%); падіння продуктивності запитів $\leq$ 20\%.
    \item НФ-005 Моніторинг (високий) --- Prometheus метрики; Grafana панелі; правила сповіщень; збереження 30 днів.
    \item НФ-006 Сумісність (середній) --- RESTful API з OpenAPI; JSON; Docker контейнери; S3-сумісний MinIO; стандартний протокол Kafka.
\end{itemize}


\section{Практична цінність}

Промислові підприємства отримують готове рішення для моніторингу обладнання на технологіях з відкритим кодом:

\begin{itemize}
    \item менше незапланованих простоїв через раннє виявлення проблем;
    \item оптимізація виробництва на основі даних;
    \item обслуговування до поломки через машинне навчання.
\end{itemize}

Інтегратори IoT можуть використовувати платформу для клієнтів без необхідності будувати все з нуля. Дослідники отримують платформу з відкритим кодом для експериментів з машинним навчанням на промислових даних з готовою інфраструктурою MLOps.

Технічні особливості платформи:

\begin{itemize}
    \item схема-на-користувача для TimescaleDB дозволяє стиснення (90-95\% економії) на відміну від захисту на рівні рядків;
    \item гібридна архітектура з нормалізованими метаданими та денормалізованими вимірюваннями забезпечує баланс між гнучкістю та швидкістю запитів;
    \item виявлення на рівні обладнання зберігає багатовимірність ознак, на відміну від аналізу кожного сенсора окремо;
    \item гібридні сповіщення поєднують пороги та моделі машинного навчання в одному механізмі.
\end{itemize}


\section{Структура роботи}

Робота складається з п'яти розділів:

\begin{itemize}
    \item Розділ 1 (Вступ) описує актуальність проблеми, мету та завдання дослідження, вимоги до системи;
    \item Розділ 2 (Огляд літератури) аналізує існуючі IoT платформи та обґрунтовує вибір технологій;
    \item Розділ 3 (Методологія) описує підхід до розробки, вибір технологій, архітектуру системи та стратегії тестування;
    \item Розділ 4 (Результати та обговорення) представляє результати тестування продуктивності, оцінку ML моделей та порівняння з існуючими рішеннями;
    \item Розділ 5 (Висновки) підсумовує досягнення цілей та окреслює напрямки подальшого розвитку.
\end{itemize}


\section{Обмеження}

Обмеження обсягу проєкту:

\begin{itemize}
    \item фокус на серверній частині та API, інтерфейс має базову функціональність;
    \item моделі машинного навчання: Isolation Forest, Random Forest, XGBoost --- глибоке навчання (LSTM, Autoencoder) залишено на майбутнє;
    \item тестування на одному сервері --- розгортання Kubernetes в перспективі.
\end{itemize}

